\section{Model I: OR-Tree Feasibility Search}
\label{sec:or}
\subsection{Search model: $A_{\lor}=(S_{\lor},T_{\lor})$}
\paragraph{Problem descriptions.}
\[
\boxed{\Prob_{\lor}=\{(a,U):a\text{ is a consistent partial assignment},\ U=X\setminus\domn(a)\}.}
\]
Here $a$ is the current partial timetable, and $U$ contains unscheduled components. A node has the lecture's form $((a,U),sol)$, where $sol\in\{\yes,\und,\no\}$.

\paragraph{Alternatives.}
Select a component $x\in U$ and define
\[
D_a(x)=\{s\in D(x):\Consistent(a\cup\{x\mapsto s\})\}.
\]
For every $s_i\in D_a(x)$, let $pr_i=(a\cup\{x\mapsto s_i\},U\setminus\{x\})$. Then the application-specific alternatives relation is
\[
\boxed{\Altern((a,U),pr_1,\ldots,pr_r),\qquad r=|D_a(x)|.}
\]
For every selected $x$, all its currently legal slot alternatives are represented; the control tries them until one succeeds or all fail.

\paragraph{States and extension rules.}
Following Lecture 11, $S_{\lor}\subseteq\Otree$ and $T_{\lor}\subseteq S_{\lor}\times S_{\lor}$. We instantiate $\Erw_{\lor}$ by the following leaf operations, which may also be performed within a larger tree:
\begin{align*}
 ((a,\varnothing),\und)&\longrightarrow((a,\varnothing),\yes)
 &&\text{if }\mathrm{Constr}(a),\\
 ((a,U),\und)&\longrightarrow((a,U),\no)
 &&\text{if no valid completion is possible},\\
 (pr,\und)&\longrightarrow(pr,\und,(pr_1,\und),\ldots,(pr_r,\und))
 &&\text{if }\Altern(pr,pr_1,\ldots,pr_r).
\end{align*}
The second case applies when the partial assignment violates a constraint or some remaining component has $D_a(x)=\varnothing$; an internal node is marked unsuccessful only after all child alternatives have failed. $\Erw_{\lor}^{*}$ allows an applicable extension inside an OR-tree while the rest of the tree is unchanged. Thus transitions apply either relation, and only one successful branch is required.

\subsection{Search process: $P_{\lor}=(A_{\lor},\Env_{\lor},K_{\lor})$}
The environment stores the scheduling instance, the current search tree, and a deterministic order for components and slots. The control selects the deepest unresolved leaf (depth-first search), then uses minimum remaining values (MRV):
\[
\boxed{x^*=\arg\min_{x\in U}|D_a(x)|.}
\]
Ties are broken by the fixed component order; feasible slots are tried in a fixed slot order. If $D_a(x^*)=\varnothing$, the leaf fails. Otherwise the control expands it with all currently legal alternatives. No forward checking is necessary for correctness.

\subsection{Search instance: $\Inst_{\lor}=(q_0,G_{\lor})$}
The initial state is $q_0=((\varnothing,X),\und)$. For a tree $q$, define success recursively:
\[
\begin{aligned}
\operatorname{Success}((pr,\yes))&=\mathsf{true},\\
\operatorname{Success}((pr,\no))&=\mathsf{false},\\
\operatorname{Success}((pr,\und,b_1,\ldots,b_r))&=\bigvee_{i=1}^{r}\operatorname{Success}(b_i).
\end{aligned}
\]
$G_{\lor}(q)=\yes$ when a successful branch exists, or when every branch has been resolved unsuccessful. These are distinct outcomes: the former returns a complete valid timetable; the latter returns FAILURE. A resource limit reached before either outcome is inconclusive.

\subsection{Implementation and operation}
The search uses a \textbf{dictionary} $a$ of assigned components and a \textbf{set} $U$ of unscheduled components. Domains are stored as lists of slot identifiers; occupancy counts and overlapping slots are cached in arrays or dictionaries for fast constraint checks. A \textbf{depth-first stack} represents the active OR-tree path. Each frame stores $(a,U)$ (or a reversible assignment trail), the MRV-selected component, untried legal slots and its status $\{\und,\yes,\no\}$. The control tries a slot, updates the counts and checks consistency; after failure it undoes the assignment and tries the next alternative. It returns the first complete valid timetable, or FAILURE only after all alternatives are exhausted.

\subsection{Worked OR-tree derivation}
Use the three-lecture instance in Section~\ref{andopt:demonstration}. Let the fixed offering order be $x<y<z$ and the OR-tree slot order be $s_2,s_1$. At the root, all three domains have size two; MRV chooses $x$. The first branch is
\[
\begin{aligned}
(\varnothing,\{x,y,z\})
&\rightarrow(\{x\mapsto s_2\},\{y,z\})\\
&\rightarrow(\{x\mapsto s_2,y\mapsto s_1\},\{z\})\\
&\rightarrow(\{x\mapsto s_2,y\mapsto s_1,z\mapsto s_2\},\varnothing).
\end{aligned}
\]
The second step has only $s_1$ legal for $y$ because $x$ and $y$ are not compatible. The final assignment satisfies capacities and all remaining hard constraints, so its leaf becomes $\yes$ and $G_{\lor}$ holds. The root's $x\mapsto s_1$ alternative remains unexplored: the OR-tree needs just one valid timetable, even though the feasible result has $\Eval=5$.
